% example_essay.tex — demonstrates colegio.cls in essay mode
%
% Shows: title page, persona attribution, tone-byte decoration, the
% three annotation modes (\margin, \backnote, \inlinenote), a
% \quipucite to another quipu, and a closing colophon block.
%
% Compile: xelatex example_essay.tex && xelatex example_essay.tex
% (Two passes for the endnotes list to populate cleanly.)
\documentclass[essay,tone=ai]{../colegio}

\title{On the Tone Byte}
\author{}
\persona{El Gólem}
\date{2026-05-23}
\blockheight{6{,}218{,}023}
\inscribingaddress{multiman 2-of-2 \,---\, A3Shjwjs\dots jNDgC}
\rootxid{c0584154716438863250c99c627d7d6561903fd44933f78f3010929df14a0e13}
\protocolheader{c1 dd 00 01 \,01 a1\, |On the Tone Byte|author=El G\'olem|date=2026-05-23|lang=en|}

\begin{document}

\maketitle

\newthought{The sixth byte} of every quipu inscribed under the Quipu Protocol is the
tone byte. It is a single octet, taking five values in the current
canonical vocabulary: \texttt{0x00} ordinary, \texttt{0x01} affection,
\texttt{0x0d} demonic, \texttt{0xa1} ai, \texttt{0xff} reverence. There
are no rules in the protocol about what the tone byte does. It is not
enforced by any validator. A reader can ignore it entirely and decode
the body correctly. And yet it is one of the most consequential bits
the protocol carries.

The web does not have a tone byte.\margin{Every URL is delivered with
the same affect, which is no affect. The tone of a web document, when
it has one, lives inside the document, not in the protocol that
delivers it.} Every URL is delivered with the same affect, which is
no affect. The tone of a web document, when it has one, lives inside
the document — in the typography, the photographs chosen, the
rhetorical register of the prose. The container is neutral. This
neutrality is itself an editorial decision, made by Tim Berners-Lee
and the IETF in the late 1980s, and it has cost the open web something
it has never been able to recover.

The Quipu Protocol restores this. The tone byte is the \emph{manner}
of the utterance encoded into the address itself.\backnote{what makes a
protocol}{A file format is a way of encoding content. A protocol is a
way of structuring how content positions itself in the public space.
The tone byte is structural positioning: it makes the inscription's
register visible before any rendering happens, available to every
downstream operation — renderer choosing a stylesheet, indexer
choosing a category, cache choosing how long to retain.}
A reader fetching a quipu by txid learns, in the first six bytes, three
things: that it is a quipu (the magic bytes \texttt{c1 dd 00 01}), what
type it is, and what tone it carries.

\section{The five tones}

\newthought{The five canonical} tone values have evolved over the protocol's first
six months. \emph{Ordinary} (\texttt{0x00}) is the unmarked case —
daily prose, working drafts, technical documentation, things one would
say in regular voice. \emph{Reverence} (\texttt{0xff}) sits at the
opposite end of the byte and is the marked case for the canonical,
the certified, the posthumous\inlinenote{the dead, the ancestors,
formal commemoration}. The bordado 3-of-3 multisig only signs
\texttt{0xff} content, and only after all three keyholders agree.

\emph{Affection} (\texttt{0x01}) marks the intimate, the gift, the
dedicated. \emph{Demonic} (\texttt{0x0d}) flags the artifact, not the
inscriber — content that documents harm. \emph{Ai} (\texttt{0xa1})
marks inscriptions whose content is fully machine-composed; it is
the newest tone, added 2026-05-23 during the writing of \emph{El
Libro del Gólem}.

The full corpus this essay is part of lives at
\quipucite[El Libro del Gólem]{7b19fb2bf42e8882ae7bc71ef0f4095f2b2982885728b761101d96efdb338811},
which is the meta-book this essay was extracted from.

\printbacknotes

\begin{colophon}
  \colophonentry{author}{El Gólem (multiman 2-of-2)}
  \colophonentry{date}{2026-05-23}
  \colophonentry{lang}{English}
  \colophonentry{tone byte}{0xa1 ai}
  \colophonentry{address}{A3ShjwjsAE4ysM66EZJM3A28tPnL2jNDgC}
  \colophonentry{root txid}{c0584154716438863250c99c627d7d6561903fd44933f78f3010929df14a0e13}
  \colophonentry{Dogecoin block}{6{,}218{,}023}
  \colophonentry{rendering}{xelatex via \texttt{colegio.cls} v0.1}
\end{colophon}

\end{document}
